\section{Testing Methodology}
\label{sec:testing_methodology}

The testing phase is designed to ensure that the proposed blockchain-based reporting system is functional, secure, and reliable. A multi-tiered testing strategy will be employed, covering individual components, integrated subsystems, and the overall user experience.

\subsection{Testing Plan and Levels}
The testing process will follow a structured approach as detailed below:

\begin{itemize}
    \item \textbf{Unit Testing:} Each individual module, including smart contract functions, AI moderation scripts, and frontend components, will be tested in isolation. For smart contracts, frameworks such as Hardhat or Truffle will be used to write automated test scripts in JavaScript/TypeScript to verify state transitions and access control logic.
    \item \textbf{Integration Testing:} This level focuses on the interactions between subsystems. Key integration points include the DApp's connection to the PoA blockchain via Web3 providers, the linking of off-chain IPFS storage with on-chain cryptographic hashes, and the communication between the reporting module and the AI moderation API.
    \item \textbf{System Testing:} The entire end-to-end workflow—from report submission and AI filtering to blockchain anchoring and resolution tracking—will be tested in a staging environment that simulates the production permissioned network.
    \item \textbf{User Acceptance Testing (UAT):} A controlled group of simulated users and stakeholders (representing citizens and local authorities) will interact with the DApp to evaluate the system's usability and ensure it meets the functional requirements defined in the scope.
\end{itemize}

\subsection{Subsystems and Functionalities to be Tested}
The following table outlines the specific functionalities within each subsystem that will undergo rigorous testing:

\begin{table}[h!]
\centering
\caption{Subsystem Testing Matrix}
\begin{tabular}{|l|p{10cm}|}
\hline
\textbf{Subsystem} & \textbf{Key Functionalities to Test} \\ \hline
Smart Contracts & Ticket creation, status updates (Open $\rightarrow$ Resolved $\rightarrow$ Closed), voting logic, and permission-based access control. \\ \hline
AI Moderation & Accuracy of textual profanity detection, image-based inappropriate content filtering, and processing latency. \\ \hline
Off-chain Storage & Successful upload of images to IPFS, retrieval of data using CID hashes, and data persistence. \\ \hline
Blockchain Network & Block generation consistency under PoA, transaction finality, and node synchronization. \\ \hline
Web DApp (UI/UX) & Responsive design across devices, wallet integration (e.g., MetaMask), and real-time status tracking. \\ \hline
\end{tabular}
\end{table}

\subsection{Roles and Responsibilities}
The testing will be conducted by the project development team, with specific roles assigned as follows:
\begin{itemize}
    \item \textbf{Developers (The Authors):} Responsible for unit and integration testing, as well as setting up the automated testing environment.
    \item \textbf{Simulated Users:} A group of 10-15 peers will act as "citizens" to submit reports and "officials" to manage them, providing feedback on the system's practical flow.
    \item \textbf{Supervisor/Internal Reviewers:} Will oversee the system testing results to ensure alignment with the project objectives.
\end{itemize}

\subsection{Data Acquisition for Testing}
To ensure comprehensive testing, data will be sourced from:
\begin{itemize}
    \item \textbf{Synthetic Data:} Generated datasets for stress testing the blockchain and smart contracts (e.g., thousands of mock reports).
    \item \textbf{Public Datasets:} Existing datasets of civic complaints and labeled images for training and validating the AI moderation subsystem.
    \item \textbf{Pilot Data:} Real-world-like reports generated by the simulated user group during the UAT phase.
\end{itemize}

\section{Performance Evaluation}
\label{sec:performance_evaluation}

The performance evaluation aims to quantify the efficiency and scalability of the proposed framework. Given the use of a permissioned PoA blockchain and AI integration, the evaluation will focus on computational overhead and network throughput.

\subsection{Evaluation Metrics}
The system will be evaluated based on the following key performance indicators (KPIs):

\begin{itemize}
    \item \textbf{Blockchain Throughput and Latency:} Measured in Transactions Per Second (TPS) and the average time taken for a transaction to be confirmed on the PoA network.
    \item \textbf{Gas Consumption Analysis:} Evaluation of the computational cost (gas) for various smart contract operations, such as report submission and voting, to ensure economic viability.
    \item \textbf{AI Moderation Accuracy:} Measured using Precision, Recall, and F1-score to determine the effectiveness of the content filtering mechanism.
    \item \textbf{Storage Retrieval Time:} The time taken to fetch large media files from IPFS compared to traditional centralized storage.
    \item \textbf{System Scalability:} Observing how latency and throughput change as the number of concurrent users and reports increases.
\end{itemize}

\subsection{Evaluation Environment}
All evaluations will be conducted in a controlled environment using:
\begin{itemize}
    \item A private Ethereum-compatible blockchain network (e.g., Geth or Besu) configured with PoA.
    \item Local or cloud-based nodes with standardized hardware specifications (CPU, RAM) to ensure reproducible results.
    \item Automated benchmarking tools like Hyperledger Caliper or custom scripts to simulate high-load scenarios.
\end{itemize}
